\chapter{Pinnaplasty (Otoplasty)}

\section*{Introduction \& Clinical Indications}
Pinnaplasty, commonly known as otoplasty, is a reconstructive and cosmetic surgical procedure primarily aimed at correcting prominent or protruding ears by reshaping or repositioning the underlying cartilaginous framework of the auricle. This intervention addresses aesthetic concerns arising from congenital deformities of the external ear, specifically the pinna, which can significantly impact a patient's self-esteem, body image, and social interactions, particularly during childhood and adolescence. The fundamental goal is to achieve a more natural-looking ear contour that lies closer to the side of the head, thereby enhancing overall facial harmony and alleviating the psychological distress often associated with prominent ears. The procedure is frequently performed bilaterally to ensure optimal symmetry, even if one ear exhibits greater prominence than the other.

\begin{itemize}
  \item Otoplasty
  \item Ear pinning surgery
  \item Prominent ear correction
  \item Setback otoplasty
  \item Auriculoplasty (a broader term encompassing various ear reshaping procedures)
  \item Cosmetic ear surgery
\end{itemize}

The fundamental surgical principle of pinnaplasty involves the precise modification of the auricle's cartilaginous framework to achieve a more aesthetically pleasing and less prominent projection from the scalp. Prominent ears typically result from one or both of two primary anatomical abnormalities: an underdeveloped or absent antihelical fold, which fails to create the natural curve that pulls the ear back towards the head, and/or an excessively deep conchal bowl, which pushes the entire ear forward. The rationale for surgical intervention is to correct these specific deformities, thereby reducing the cephaloauricular angle (the angle between the mastoid and the ear's posterior surface) to a more normal range of 20-30 degrees.

Correction of an underdeveloped antihelical fold is achieved by creating or enhancing this fold. Techniques involve either weakening the cartilage through scoring or abrasion on its anterior surface, allowing it to bend into the desired shape, or using permanent mattress sutures (e.g., Mustard\'e sutures) placed through the cartilage to physically fold and maintain the antihelix. For an overly deep conchal bowl, the surgeon may excise a crescent-shaped piece of cartilage from the concha to reduce its volume, or employ permanent sutures (e.g., Furnas sutures) to pull the conchal cartilage closer to the mastoid fascia, effectively setting the ear back. The overarching technical goals are to achieve a natural-looking ear contour without an "over-pinned" appearance, ensure optimal symmetry between both ears, minimize visible scarring by placing incisions discreetly, and maintain the structural integrity and long-term stability of the ear's new position. The specific technique or combination of techniques is meticulously chosen based on the individual patient's unique anatomical abnormalities identified during the preoperative assessment.

The history of otoplasty dates back to the late 19th century, with early surgical attempts primarily focused on simple excision of skin and cartilage to reduce ear prominence. One of the earliest documented surgical corrections for prominent ears was performed by Dr. Edward T. Ely in 1881, who described a technique involving skin excision and cartilage suturing. However, these initial methods often yielded inconsistent results, frequently leading to unnatural-looking ears, sharp edges, or a high rate of recurrence due to the inherent memory of the cartilage.

Significant advancements in understanding the underlying anatomical deformities and developing more refined techniques emerged in the early 20th century. In 1910, Dr. Albert Luckett recognized the critical role of the antihelical fold in ear prominence and proposed a technique to create it by excising a portion of the conchal cartilage. The mid-20th century marked a pivotal era with the introduction of sophisticated cartilage-sparing and cartilage-reshaping methods. In 1963, Dr. John Mustard\'e revolutionized otoplasty with his seminal technique, which utilized permanent mattress sutures to create and maintain the antihelical fold without excising cartilage, relying instead on the cartilage's inherent memory and the tension of the sutures. Shortly thereafter, in 1968, Dr. David Furnas described a complementary technique to reduce conchal prominence by suturing the conchal cartilage directly to the mastoid periosteum. These pioneering contributions laid the foundation for modern otoplasty, which frequently combines elements of both cartilage-suturing and cartilage-scoring methods to achieve optimal, natural, and long-lasting results, tailored to the specific anatomical needs of each patient.

Pinnaplasty is primarily indicated for aesthetic and psychosocial reasons, addressing concerns related to the appearance of the ears. Specific clinical indications include:

\begin{itemize}
  \item Correction of prominent or protruding ears that extend excessively from the side of the head, often due to an increased cephaloauricular angle.
  \item Improvement of ear asymmetry, where one ear is significantly more prominent or differently shaped than the other, leading to an unbalanced facial appearance.
  \item Reshaping of ears with an underdeveloped, absent, or poorly defined antihelical fold, which results in a flat or "unfolded" ear appearance.
  \item Reduction of an overly deep conchal bowl (conchal hypertrophy), which pushes the entire ear forward and contributes significantly to ear protrusion.
  \item Alleviation of psychological distress, embarrassment, low self-esteem, or social anxiety experienced by individuals, particularly children and adolescents, due to their prominent ear appearance.
  \item Prevention or mitigation of teasing and bullying that children often experience in school settings because of their ear appearance.
  \item Reconstruction of ears following trauma or other congenital deformities beyond simple prominence, although these cases may involve more complex reconstructive techniques.
  \item Enhancement of overall facial balance and harmony by bringing the ears into better proportion with other facial features.
\item The procedure is typically considered once the ear cartilage is sufficiently developed, usually around 5-7 years of age, to ensure stable results and avoid interference with ear growth.
\end{itemize}

Pinnaplasty is overwhelmingly considered a primary, elective cosmetic surgical procedure. It is typically the first and often the only surgical intervention undertaken to correct prominent ears, chosen by patients or their parents primarily for aesthetic improvement and psychosocial well-being. Prominent ears are a congenital condition that does not typically require prior medical or non-surgical treatment, making otoplasty a first-line solution for definitive correction. The decision to proceed is usually driven by the patient's or family's desire to address the cosmetic appearance and mitigate potential psychological impacts.

In rare instances, a revision otoplasty might be considered a secondary or salvage procedure. This occurs if the initial surgery yields an unsatisfactory cosmetic result, such as recurrence of prominence, significant asymmetry, an unnatural contour (e.g., "telephone ear" deformity), or complications like suture extrusion. In such cases, the secondary procedure aims to refine, correct, or salvage the outcome of the primary surgery, often requiring a more complex approach due to altered anatomy and scar tissue. However, for the vast majority of patients seeking correction of prominent ears, pinnaplasty serves as the initial and definitive treatment.

\section*{Relevant Surgical Anatomy}
The external ear, or auricle (pinna), is a complex three-dimensional structure primarily composed of elastic cartilage covered by thin skin. Its intricate folds and depressions are crucial for both aesthetics and sound collection. Key cartilaginous landmarks include the helix (the outer rim), the antihelix (an inner fold parallel to the helix, which normally creates a natural curve that pulls the ear back), the scaphoid fossa, the triangular fossa, the concha (a deep central bowl divided into the cymba conchae superiorly and cavum conchae inferiorly), the tragus, and the antitragus. The earlobe, or lobule, is unique as it lacks cartilage, consisting instead of fat and fibrous tissue. The skin covering the auricle is tightly adherent on the anterior surface and looser on the posterior surface, where the surgical incision is typically made.

The blood supply to the auricle is robust, primarily derived from branches of the external carotid artery. The posterior auricular artery, a branch of the external carotid, supplies the posterior surface of the ear and contributes to the anterior surface. The superficial temporal artery, another external carotid branch, supplies the anterior and superior aspects. Venous drainage largely parallels the arterial supply, with the posterior auricular vein draining into the external jugular vein and the superficial temporal vein draining into the retromandibular vein. Sensory innervation is complex and critical to preserve. The greater auricular nerve (C2, C3 from the cervical plexus) provides sensation to the posterior surface, the helix, antihelix, and lobule. The auriculotemporal nerve (a branch of the mandibular nerve, V3) innervates the tragus, crus of the helix, and anterior superior helix. The lesser occipital nerve (C2) supplies a small portion of the superior helix. Motor innervation to the vestigial auricular muscles is provided by branches of the facial nerve (CN VII), which must be carefully avoided during dissection, particularly near the mastoid process. Lymphatic drainage occurs via the preauricular, mastoid (posterior auricular), and deep cervical lymph nodes. Surgical landmarks include the postauricular sulcus for incision placement and the mastoid fascia for conchal setback sutures.

\section*{Preoperative Workup \& Preparation}
Thorough preoperative workup and preparation are paramount to ensure patient safety, optimize surgical outcomes, and manage expectations effectively. This process typically involves several key steps:

\begin{itemize}
  \item \textbf{Comprehensive Medical Evaluation:} A detailed medical history is obtained, and a physical examination is performed to assess the patient's overall health status. For children, a pediatrician's clearance confirming fitness for general anesthesia is often required. Any pre-existing medical conditions, allergies, or previous surgeries are thoroughly reviewed.
  \item \textbf{Laboratory Tests:} Routine preoperative blood tests, such as a complete blood count (CBC), basic metabolic panel, and coagulation profile (PT/INR, PTT), may be ordered to screen for underlying medical conditions, anemia, or bleeding disorders.
  \item \textbf{Anesthesia Clearance:} An anesthesiologist conducts a pre-anesthetic evaluation to determine the most appropriate type of anesthesia (general vs. local with sedation) and to ensure the patient is medically fit for the planned procedure.
  \item \textbf{Medication Review and Adjustment:} Patients must provide a complete list of all prescription medications, over-the-counter drugs, herbal remedies, and supplements. Blood-thinning medications (e.g., aspirin, NSAIDs like ibuprofen, warfarin, clopidogrel) are typically discontinued 1-2 weeks prior to surgery to minimize the risk of bleeding, under the strict guidance of the prescribing physician.
  \item \textbf{Smoking Cessation:} Patients who smoke are strongly advised to cease smoking for at least 4-6 weeks before surgery, as nicotine can impair wound healing, increase the risk of infection, and compromise blood supply to the tissues.
  \item \textbf{NPO Status:} Patients are instructed to refrain from eating or drinking (nil per os - NPO) for a specified period, usually 6-8 hours, before surgery to prevent aspiration during anesthesia induction.
  \item \textbf{Prophylactic Antibiotics:} A single dose of intravenous prophylactic antibiotics (e.g., cefazolin) is often administered preoperatively to reduce the risk of surgical site infection, particularly chondritis.
  \item \textbf{Informed Consent:} The surgeon conducts a detailed discussion with the patient (or legal guardian) regarding the procedure, including its anticipated benefits, potential risks, possible complications, expected outcomes, and alternative treatments. A comprehensive informed consent form must be signed.
  \item \textbf{Preoperative Photography:} Standardized preoperative photographs are taken from multiple angles for medical record documentation, surgical planning, and objective assessment of postoperative results.
  \item \textbf{Hygiene:} Patients are advised to shower with an antiseptic soap the night before or morning of surgery.
\end{itemize}

Careful consideration of contraindications and precautions is paramount to ensure patient safety and achieve optimal surgical outcomes in pinnaplasty.

\textbf{Absolute Contraindications:}
\begin{itemize}
  \item Active infection in or around the ear: Any localized or systemic infection, especially cellulitis or otitis externa, necessitates postponement of surgery until the infection is completely resolved to prevent spread and severe complications like chondritis.
  \item Uncontrolled bleeding disorders: Patients with severe coagulopathies or those on chronic anticoagulation that cannot be safely interrupted are at a significantly elevated risk for hematoma formation, which can compromise the surgical result and lead to skin necrosis.
  \item Unrealistic patient expectations: Patients exhibiting body dysmorphic disorder or having expectations that cannot be realistically met surgically are generally considered poor candidates, as they are likely to be dissatisfied with the outcome regardless of technical success.
  \item Severe underlying medical conditions: Any unstable or severe systemic medical condition (e.g., uncontrolled heart disease, severe respiratory compromise) that precludes safe administration of anesthesia or significantly increases surgical risk.
\end{itemize}

\textbf{Relative Contraindications and Precautions:}
\begin{itemize}
  \item Age: While otoplasty can be performed from around age 5-6, when ear cartilage is sufficiently developed and stable, performing it too early (before age 5) may risk affecting natural ear growth, increasing the likelihood of recurrence, or making the child less cooperative.
  \item Minor acute illnesses: A common cold, flu, or other minor acute illness may necessitate postponing surgery to avoid respiratory complications under anesthesia and ensure the patient is in optimal health for recovery.
  \item Certain medications: Anticoagulants or antiplatelet agents require careful management, often involving temporary cessation before surgery, which must be balanced against the patient's risk of thrombotic events.
  \item History of keloid or hypertrophic scarring: Patients with a known predisposition to abnormal scarring require careful counseling regarding scar risks and may benefit from adjunctive treatments (e.g., steroid injections, silicone sheeting) postoperatively.
  \item Psychological instability: Patients with significant psychological conditions, beyond simple self-consciousness, may require prior psychological evaluation and treatment to ensure they are appropriate candidates and have realistic expectations.
  \item Active autoimmune diseases or immunosuppression: These conditions can impair wound healing, increase the risk of infection, and may necessitate specific perioperative management strategies.
  \item Previous ear surgery: Prior otoplasty or ear trauma can alter anatomy and tissue planes, making revision surgery more challenging and potentially increasing complication rates.
\end{itemize}

\section*{Surgical Technique \& Procedure}
Pinnaplasty is typically performed on an outpatient basis, with the specific technique tailored to the patient's individual anatomy and the surgeon's preference. For children, general anesthesia is almost always employed to ensure complete immobility and comfort. Adults may opt for local anesthesia with intravenous sedation, or general anesthesia.

The patient is positioned supine, and the head is turned to expose the ear. The surgical area is meticulously prepped with an antiseptic solution and draped. The procedure generally follows these key steps:

\begin{enumerate}
  \item \textbf{Anesthesia:} Administered as general anesthesia for pediatric patients or local anesthesia with intravenous sedation for adults. Local anesthetic with epinephrine is infiltrated into the posterior surface of the ear and around the conchal bowl to provide intraoperative analgesia and vasoconstriction.
  \item \textbf{Incision:} A curvilinear incision is typically made on the posterior surface of the ear, within the natural postauricular crease where the ear meets the head. This strategic placement ensures that any resulting scar is well-hidden and inconspicuous.
  \item \textbf{Cartilage Exposure \& Dissection:} The skin and subcutaneous tissue are carefully lifted from the underlying ear cartilage, exposing the posterior aspect of the auricle. Minimal undermining is performed to preserve blood supply.
  \item \textbf{Cartilage Reshaping (Antihelical Fold Creation):} The surgeon assesses the cartilage. If an underdeveloped antihelical fold is the primary issue, techniques may include:
    \begin{itemize}
      \item \textit{Cartilage Scoring/Weakening:} Small, superficial incisions or abrasion are made on the anterior surface of the cartilage (accessed via the posterior incision) to weaken its elastic memory, allowing it to bend more easily into the desired antihelical fold.
      \item \textit{Suture Techniques (e.g., Mustard\'e):} Permanent, non-absorbable mattress sutures are placed through the cartilage from the posterior aspect to create and maintain the antihelical fold, without excising cartilage. These sutures are tied to achieve the desired contour.
    \end{itemize}
  \item \textbf{Conchal Setback (Conchal Reduction):} If an excessively deep conchal bowl contributes to prominence, the surgeon may:
    \begin{itemize}
      \item \textit{Cartilage Excision:} A small, crescent-shaped piece of cartilage may be excised from the concha to reduce its depth and projection.
      \item \textit{Suture Techniques (e.g., Furnas):} Permanent sutures are placed from the conchal cartilage to the mastoid fascia (the tough tissue covering the bone behind the ear) to pull the concha closer to the head, effectively reducing its prominence.
    \end{itemize}
  \item \textbf{Symmetry Check \& Refinement:} The surgeon meticulously checks the contour, projection, and symmetry of both ears against each other and the patient's facial features to ensure an optimal and natural appearance. Adjustments are made as necessary.
  \item \textbf{Closure:} The skin incision behind the ear is closed with fine, absorbable or non-absorbable sutures. Care is taken to achieve a neat closure that minimizes tension and optimizes scar healing.
  \item \textbf{Dressing Application:} A soft, bulky head dressing or bandage is applied to protect the ears, maintain their new position, and minimize swelling and hematoma formation. Drains are typically not used in routine pinnaplasty.
\end{enumerate}

\section*{Complications \& Risks}
As with any surgical procedure, pinnaplasty carries potential risks and complications, which patients must be thoroughly informed about during the consent process. These can be broadly categorized into general surgical risks and procedure-specific risks.

\begin{itemize}
  \item \textbf{General Surgical Risks:}
    \begin{itemize}
      \item \textit{Bleeding and Hematoma:} Accumulation of blood under the skin or within the cartilage, which can cause pain, swelling, and potentially compromise the surgical result. A significant hematoma may require surgical drainage.
      \item \textit{Infection:} Although rare, infection of the surgical site or, more severely, the cartilage (chondritis) can occur. Chondritis is difficult to treat and can lead to cartilage destruction and significant deformity. It typically requires aggressive antibiotic therapy and sometimes surgical debridement.
      \item \textit{Anesthesia Complications:} Risks associated with general or local anesthesia, including nausea, vomiting, allergic reactions, respiratory depression, or more serious cardiovascular events.
      \item \textit{Pain:} Postoperative pain is expected but usually manageable with oral analgesics. Persistent or severe pain should be evaluated for complications.
      \item \textit{Scarring:} While incisions are typically hidden behind the ear, abnormal scarring such as hypertrophic scars (raised, red scars confined to the incision) or keloids (scars that grow beyond the incision boundaries) can occur, especially in predisposed individuals.
    \end{itemize}
  \item \textbf{Procedure-Specific Risks:}
    \begin{itemize}
      \item \textit{Asymmetry:} Despite meticulous surgical technique, perfect symmetry between the two ears is challenging to achieve, and minor differences in shape, size, or projection may persist or develop.
      \item \textit{Recurrence of Prominence:} The ear cartilage has a "memory" and may gradually return to its original prominent position, especially if sutures fail, cartilage weakening is insufficient, or postoperative care is not followed. This may necessitate revision surgery.
      \item \textit{Unsatisfactory Cosmetic Result:} The ears may appear overcorrected (too flat, "pinned-back" or "telephone ear" deformity where the middle portion is too close to the head while the top and bottom protrude), undercorrected, or have an unnatural, irregular, or sharp contour.
      \item \textit{Numbness or Altered Sensation:} Temporary or, rarely, permanent numbness or altered sensation in parts of the ear due to irritation or damage to sensory nerves (e.g., greater auricular nerve) during dissection.
      \item \textit{Suture Extrusion:} Permanent sutures used to reshape the cartilage may occasionally become visible, palpable, or extrude through the skin, requiring removal.
      \item \textit{Skin Necrosis:} A rare but serious complication where skin tissue dies due to compromised blood supply, potentially requiring wound care, debridement, or further reconstructive intervention.
      \item \textit{Chondritis:} As mentioned, infection of the cartilage is a severe complication that can lead to significant deformity and requires prompt, aggressive treatment.
      \item \textit{Allergic Reaction:} To suture materials, dressings, or topical antibiotics.
      \item \textit{Pressure Sores:} Can occur under a tight head dressing if not properly applied or monitored.
    \end{itemize}
\end{itemize}

\section*{Postoperative Care \& Recovery}
Immediately following pinnaplasty, the patient will be transferred to the Post-Anesthesia Care Unit (PACU) for close monitoring as they recover from anesthesia. A bulky head dressing or bandage will be securely in place to protect the ears, minimize swelling, and help maintain the newly established ear position. Pain medication will be administered as needed to manage discomfort. Most patients are discharged home on the same day, provided they are stable, alert, and have adequate pain control, along with detailed instructions for wound care and activity restrictions.

During the first 24-48 hours, some discomfort, mild to moderate pain, and swelling are expected. Oral analgesics are typically sufficient for pain management. The head dressing usually remains in place for approximately 5-7 days. Patients are strictly advised to keep the dressing dry and avoid any pressure or trauma to the ears. After the initial dressing is removed by the surgeon, a lighter, soft elastic headband (similar to a sports headband) is typically recommended. This headband should be worn continuously for the first 1-2 weeks (except for showering) and then nightly for an additional 4-6 weeks to provide support during the critical healing phase and prevent accidental injury or displacement of the ears while sleeping. Most children can return to school within a week, while adults may return to work sooner, depending on their occupation. Strenuous activities, contact sports, and activities that could put pressure on the ears must be avoided for at least 6-8 weeks. Swelling and bruising will gradually subside over several weeks, and the final, refined ear contour will become more apparent as healing progresses over several months.

During pinnaplasty, the surgeon's primary findings relate to the specific anatomical deformities contributing to ear prominence. These typically include an underdeveloped or absent antihelical fold, which results in a flat or unfolded appearance of the ear's upper pole, and/or an excessively deep conchal bowl (conchal hypertrophy) that pushes the entire auricle away from the mastoid. The surgeon meticulously assesses the degree of these deformities and notes any asymmetries between the two ears. Other findings might include minor irregularities in the cartilage or skin.

Pathological findings are generally not relevant for routine pinnaplasty, as the procedure primarily involves reshaping normal cartilage rather than excising diseased tissue. Tissue specimens are rarely sent for pathological examination unless an unexpected lesion, such as a cyst, tumor, or suspicious skin abnormality, is encountered during the dissection. In such rare instances, the pathology report would describe the nature of the excised tissue. For the vast majority of otoplasty cases, the "pathology" is purely anatomical, referring to the structural variations of the cartilage that lead to the prominent ear appearance, and the surgical documentation serves as the primary record of these findings and their correction.

A normal and successful postoperative course following pinnaplasty is characterized by a predictable progression of healing and gradual resolution of initial symptoms, leading to the desired aesthetic outcome. Immediately after surgery, patients will experience mild to moderate pain, swelling, and bruising, which are well-managed with prescribed oral analgesics and cold compresses. The initial bulky head dressing provides protection and support. Once the dressing is removed, typically within a week, the ears will appear closer to the head, though some residual swelling will obscure the final contour.

Over the subsequent weeks, swelling and bruising will progressively diminish. The use of a protective headband, initially continuous and then nightly, is crucial to maintain the new ear position and prevent trauma. Patients will gradually regain normal sensation in the ears, although some temporary numbness or altered sensation is common. The incision behind the ear will heal, forming a fine, inconspicuous scar within the natural crease. By 6-8 weeks, most of the swelling will have resolved, and the ears will have settled into their new, less prominent position, with a more defined antihelical fold and reduced conchal depth. The final aesthetic result, including complete scar maturation, typically takes 6-12 months to fully manifest, at which point the patient should experience improved facial harmony and resolution of the psychosocial distress that prompted the surgery.

Meticulous postoperative care and regular follow-up appointments are essential to monitor healing, address any concerns, and ensure optimal long-term results after pinnaplasty.

\begin{itemize}
  \item \textbf{Initial Postoperative Visits:} The first follow-up appointment is typically scheduled within 5-7 days for removal of the initial bulky head dressing, wound inspection, and removal of any non-dissolvable skin sutures. Subsequent visits are usually scheduled at 2-4 weeks to assess early healing, monitor for complications, and reinforce instructions.
  \item \textbf{Headband Use Guidance:} Detailed instructions are provided on the duration and frequency of wearing the protective headband (e.g., continuously for 1-2 weeks, then nightly for 4-6 weeks) to support the ears and prevent recurrence or injury.
  \item \textbf{Activity Resumption:} Patients are advised on a gradual return to normal activities. Strenuous exercise and activities that could put pressure on the ears or risk trauma (e.g., contact sports) must be avoided for at least 6-8 weeks.
  \item \textbf{Wound Care:} Instructions for keeping the incision clean and dry, and for applying any prescribed topical ointments, are provided.
  \item \textbf{Long-term Follow-up:} Further visits may be scheduled at 3-6 months and one year post-surgery to assess the final aesthetic outcome, evaluate scar maturation, and address any late concerns or potential complications.
  \item \textbf{Warning Signs:} Patients are educated to contact their surgeon immediately if they experience any signs of infection (e.g., fever, increasing redness, warmth, severe pain, pus drainage), excessive swelling, significant bleeding, sudden changes in ear sensation, or any displacement of the ear's position.
  \item \textbf{Revision Surgery Discussion:} In the rare event of an unsatisfactory cosmetic result, such as significant asymmetry or recurrence of prominence, revision surgery may be discussed after a sufficient healing period (typically 6-12 months) has elapsed.
\end{itemize}
Adherence to these postoperative guidelines is crucial for achieving a successful and lasting outcome.

\section*{Outcomes \& Prognosis}
Prominent ears represent one of the most common congenital deformities of the head and neck, affecting approximately 5\% of the general population. The condition often exhibits a familial pattern, suggesting a significant genetic predisposition, although the exact mode of inheritance can be variable and multifactorial. There is no significant predilection for sex, with prominent ears occurring almost equally in males and females across various ethnic groups. The primary indication for pinnaplasty is cosmetic correction of ear protrusion, driven predominantly by the profound psychosocial impact on affected individuals, particularly children who frequently experience teasing, bullying, and diminished self-esteem, which can negatively affect their social development and academic performance. The procedure is most frequently performed in children between the ages of 5 and 14 years, as the ear cartilage is nearly fully developed by age 5-6, making it structurally stable enough for surgical reshaping, and early intervention can effectively mitigate psychological distress during formative school years. While less common, adults also seek pinnaplasty for similar aesthetic and psychological reasons. The morbidity associated with pinnaplasty is generally low, with minor complications such as hematoma, infection, and temporary numbness occurring in a small percentage of cases. Mortality is exceedingly rare and is typically associated with general anesthesia risks rather than the procedure itself.

The outcomes of pinnaplasty are generally excellent, with high rates of patient satisfaction and significant improvement in ear aesthetics and psychosocial well-being. Typical success rates, defined by satisfactory cosmetic results and patient contentment, range from 85\% to 95\% across various studies and surgical techniques. Functional outcomes are not a primary concern, as the procedure is purely cosmetic, but patients consistently report improved self-confidence, reduced social anxiety, and a better quality of life following successful surgery. The prognosis for long-term stability of the correction is good, especially with modern techniques that effectively address both antihelical fold deficiency and conchal excess, often combining cartilage-scoring and suture-based methods. Recurrence of prominence, while possible due to cartilage memory or suture failure, is relatively low, estimated at 5\% to 10\%, and may necessitate revision surgery.

Prognostic factors influencing outcomes include the surgeon's experience and expertise, the specific technique or combination of techniques employed, the severity and type of the initial deformity, and crucial patient compliance with postoperative care instructions, particularly the consistent use of the protective headband. While minor asymmetries can persist, significant complications leading to permanent disfigurement are rare. Overall, pinnaplasty offers a predictable and highly effective solution for prominent ears, providing lasting aesthetic and psychological benefits that significantly enhance the patient's self-perception and social interactions.

\section*{References}
\begin{enumerate}
  \item MedlinePlus. Otoplasty. U.S. National Library of Medicine; 2024. Available from: \url{https://medlineplus.gov/earsurgery.html}
  \item Thorne, C. H. (Ed.). (2017). \textit{Grabb and Smith's Plastic Surgery} (8th ed.). Wolters Kluwer.
  \item American Society of Plastic Surgeons. Otoplasty (Ear Surgery). Available from: \url{https://www.plasticsurgery.org/cosmetic-procedures/ear-surgery}
  \item Janis, J. E. (Ed.). (2014). \textit{Essentials of Plastic Surgery} (2nd ed.). Thieme.
\end{enumerate}

\clearpage
